% OVERLEAF : New Project > Blank Project > coller dans main.tex > Recompile (pdfLaTeX)
% Offre : Ingénieur Développement Logiciel C++ — Applications 3D & Intelligence Artificielle (F/H) — réf. 550662
\documentclass[10pt,a4paper]{article}
\usepackage[margin=0.95cm]{geometry}
\usepackage[T1]{fontenc}
\usepackage[utf8]{inputenc}
\usepackage{hyperref,enumitem,titlesec}
\hypersetup{hidelinks}
\pagestyle{empty}
\setlength{\parindent}{0pt}
\titleformat{\section}{\bfseries\large\uppercase}{}{0em}{}[\vspace{0.5pt}\titlerule\vspace{2pt}]
\titlespacing*{\section}{0pt}{3pt}{1.5pt}
\setlist[itemize]{leftmargin=*,itemsep=0.2pt,topsep=0.4pt,parsep=0pt}

\begin{document}
\small

\begin{center}
{\LARGE\bfseries Ghislain THEUBO}\\[1pt]
{\large Ing\'enieur D\'eveloppement Logiciel --- C++, Applications 3D \& Intelligence Artificielle}\\[2pt]
+33 6 82 86 62 79 $\cdot$
\href{mailto:bergerghislaintheubo@gmail.com}{bergerghislaintheubo@gmail.com} $\cdot$
\href{https://www.linkedin.com/in/berger-ghislain}{LinkedIn} $\cdot$
\href{https://github.com/Bergerghislain}{GitHub}\\[1pt]
Paris / \^Ile-de-France --- mobile V\'elizy $\cdot$ Dispo imm\'ediate $\cdot$ FR courant / EN bilingue $\cdot$ R\'ef. 550662
\end{center}

\section{Profil}
Ing\'enieur Bac+5 (Grenoble INP -- Ensimag, g\'enie logiciel) avec une \textbf{exp\'erience R\&D sur la plateforme 3DEXPERIENCE} (Dassault Syst\`emes) et une pratique concr\`ete du \textbf{C/C++} sur syst\`emes r\'eels (INRIA : temps r\'eel, flux capteurs, latence). Socle \textbf{algorithmique et math\'ematique} solide (g\'eom\'etrie, maillages, arbres spatiaux, graphes) et exp\'erience \textbf{Python / IA} en production. Cherche pr\'ecis\'ement le croisement \textbf{applications 3D $\times$ IA} : repr\'esentation g\'eom\'etrique, performance temps r\'eel et int\'egration de mod\`eles au cœur d'un logiciel industriel.

\section{Exp\'eriences professionnelles}
\textbf{Dassault Syst\`emes --- 3DEXPERIENCE / 3DSwym R\&D} \hfill V\'elizy-Villacoublay\\
\textit{Ing\'enieur logiciel (R\&D) --- refonte compl\`ete d'un module de la plateforme} \hfill Mars 2025 -- Sept. 2025
\begin{itemize}
  \item Refonte end-to-end d'un module \textbf{3DEXPERIENCE} : analyse de l'existant, sp\'ecifications techniques, architecture, d\'eveloppement, tests, livraison dans une \textbf{base de code industrielle large} (contraintes de compatibilit\'e et de maintenabilit\'e long terme).
  \item \textbf{Optimisation de performance $\sim$35\,\%} sur les chemins critiques (algorithmique, structures de donn\'ees, r\'eduction des acc\`es co\^uteux) --- d\'emarche directement transposable aux traitements g\'eom\'etriques et au rendu.
  \item Qualit\'e : \textbf{TDD}, couverture $\sim$\textbf{85\,\%}, revues de code, CI, documentation technique maintenue \`a jour.
  \item \'Equipe R\&D $\sim$20 personnes : ateliers Agile / MoSCoW avec les \textbf{product Champions}, arbitrages techniques, communication en anglais au quotidien.
\end{itemize}

\textbf{INRIA Grenoble --- Humanlab (projet Froll'N'Roll)} \hfill Montbonnot-Saint-Martin\\
\textit{Stagiaire recherche --- logiciel embarqu\'e temps r\'eel (C/C++, Python)} \hfill Mai 2024 -- Ao\^ut 2024
\begin{itemize}
  \item D\'eveloppement de modules \textbf{C / FreeRTOS} pilotant un syst\`eme physique : traitement de \textbf{flux capteurs} sous contraintes de latence et de fiabilit\'e.
  \item Refonte des protocoles et des structures de donn\'ees : \textbf{$\sim$40\,\%} de r\'eactivit\'e gagn\'ee (mesur\'ee) --- profiling, suppression de copies, ordonnancement des t\^aches.
  \item Outillage et tests en \textbf{Python}, validation it\'erative avec les utilisateurs finaux, code document\'e et ouvert ; collaboration pluridisciplinaire (recherche, mat\'eriel, usage).
\end{itemize}

\textbf{Gustave Roussy \& CentraleSup\'elec MICS} \hfill Paris / Saclay\\
\textit{Software Engineer (CDD) --- services de donn\'ees, Python \& IA appliqu\'ee} \hfill Jan. 2026 -- Sept. 2026
\begin{itemize}
  \item Conception et d\'eveloppement de services \textbf{Python} (APIs, traitements, validation, contr\^ole d'acc\`es) pour des utilisateurs experts (recherche / clinique).
  \item Mod\'elisation \textbf{SQL / PostgreSQL} (sch\'emas, index, transactions), optimisation de requ\^etes (\textbf{$\sim$30\,\%} sur les lectures critiques) et interfaces \textbf{TypeScript} de visualisation et de pilotage.
  \item Industrialisation : Docker, CI (GitHub Actions), tests, journalisation structur\'ee, diagnostic d'incidents, documentation et maintien en condition op\'erationnelle.
  \item Usage quotidien d'outils \textbf{IA} (Cursor, Claude, Copilot) pour impl\'ementation, refactoring et compr\'ehension de code, avec \textbf{relecture critique syst\'ematique} avant int\'egration.
\end{itemize}

\textbf{Shift Technology --- Platform Software Engineering} \hfill Paris, Sept. 2026 -- Nov. 2026\\
Services et APIs haute performance (C\#/.NET) sur plateforme SaaS multi-\'equipes : revues, livraisons incr\'ementales, ownership production.

\textbf{LIG --- Laboratoire d'Informatique de Grenoble} \hfill Grenoble, F\'evr. -- Mars 2025\\
Stage cloud : automatisation Linux, migration d'environnements OpenStack / Grid'5000, \textbf{mesures de performance} et reproductibilit\'e des exp\'eriences.

\section{Projets personnels --- C++, g\'eom\'etrie \& IA}
\begin{itemize}
  \item \textbf{Biblioth\`eque de maillages polyh\'edriques} (C++, templates, Catch2) : structures g\'eom\'etriques, requ\^etes spatiales par \textbf{KD-Tree}, algorithmes de graphes, tests unitaires et benchmarks.
  \item \textbf{Pipeline de transformation haute volum\'etrie} (C++, multithreading) : pools de workers, back-pressure, m\'etriques ; comparaison s\'equentiel / parall\`ele et analyse des goulots d'\'etranglement.
  \item \textbf{IA appliqu\'ee} (Python) : inf\'erence et \'evaluation de mod\`eles, int\'egration par API ; mont\'ee en comp\'etence active sur l'\textbf{inf\'erence embarqu\'ee en C++} et les APIs graphiques (OpenGL / rendu temps r\'eel).
\end{itemize}

\section{Comp\'etences techniques}
\textbf{Langages :} \textbf{C++} (POO, STL, templates), C, Python, Java, C\#, TypeScript, SQL, Bash\\
\textbf{3D \& algorithmique :} g\'eom\'etrie et maillages, KD-Tree, graphes, alg\`ebre lin\'eaire, optimisation de performance, temps r\'eel ; OpenGL / rendu (mont\'ee en comp\'etence)\\
\textbf{IA :} Python (NumPy, scikit-learn), int\'egration de mod\`eles par API, \'evaluation ; d\'eveloppement assist\'e par IA avec revue critique\\
\textbf{Qualit\'e \& outils :} TDD, Catch2 / pytest, CI/CD, Git, CMake, Docker, Linux, profiling \& d\'ebogage, documentation technique

\section{Formation}
\textbf{Grenoble INP -- Ensimag} --- Dipl\^ome d'ing\'enieur, Informatique \& G\'enie logiciel \hfill 2023 -- 2025\\
\textbf{\'Ecole Nationale Sup\'erieure Polytechnique de Yaound\'e} --- Master G\'enie informatique (Top 12\,\%) \hfill 2019 -- 2023\\
Certification Test-Driven Development (2025) $\cdot$ \textbf{Anglais} bilingue $\cdot$ \textbf{Fran\c{c}ais} courant $\cdot$ \textbf{Espagnol} interm\'ediaire

\end{document}
